\subsection{Guided search with \enumo}
\label{subsec:eval-enumo}

To evaluate \enumo's guided search,
  we conducted the following experiments on a
  64-bit Linux machine with 32 GB RAM,
  running Ubuntu 22.04.2 LTS.

\begin{table}[h]
\resizebox{\textwidth}{!}{%
\begin{tabular}{lcllcc}
Domain   & \enumo LOC  & \# \enumo (Time)   & \# \ruler (Time)  & \enumo $\rightarrow$ \ruler (Time) & \ruler $\rightarrow$ \enumo (Time)  \\ \cline{1-6}
bool & 41 & 64 (0.35)   & 51 (0.05) & 100\% (0.01), 100\% (0.01) & 87.5\% (5.29), 96.9\% (0.01) \\
bv4  & 19 & 180 (7.13)  & 84 (0.96) & 100\% (0.17), 100\% (0.03) & 38.3\% (3.67), 41.1\% (4.32) \\
bv32 & 18 & 120 (48.78) & 78 (13.1) & 100\% (0.15), 100\% (0.01)  & 58.3\% (1.41), 60.0\% (2.08) \\
rational & 57 & 123 (6.34) & 113 (97.9) & 97.3\% (0.73), 100\% (0.09) & 52.0\% (18.66), 58.5\% (22.24) \\
\end{tabular}%
}
\caption{Results comparing \enumo to \ruler.
  $ R_1 ~ \rightarrow ~ R_2$ indicates using $R_1$ to derive
  $R_2$ rules.
  We report both \lhs and \lhsandrhs derivability (in that order), separated by commas.
  The numbers in parentheses are times in seconds.
}
\label{table:oopsla}
\end{table}

\subsubsection{Comparing rulesets with prior work}
\label{subsubsec:e-vs-r}
\ruler~\citep{ruler} is a state-of-the-art
  tool for automatically
  synthesizing rewrite rules targeted towards \eqsat-driven
  systems.
We compare the rules generated by \enumo
  against those generated by \ruler\footnote{bool, bv4, and bv32
  rules are taken from the \ruler artifact. The artifact did not contain
  rational rules, so those were generated following the instructions in
  the \ruler repository on GitHub.},
  finding that
  rulesets from small \enumo programs outperform
  those from \ruler,
  which uses a hard-coded rule-finding
  strategy.
We wrote \enumo programs for
  each of the domains showcased in \ruler:
  \texttt{bool}, \texttt{bv4}, \texttt{bv32}, and \texttt{rational}.
These programs call \T{recursive_rules},
an \enumo-provided utility function (\autoref{fig:rec-rules}).
From a user-provided grammar $\mathcal{G}$ that specifies literal terms,
  unary operators, and binary operators, \T{recursive_rules} builds
  workloads of increasing size, then finds and validates rules from
  the workloads, using rules it finds along the way to avoid redundancy
  in the final ruleset.
This algorithm
 replicates \ruler's core loop in just a few lines of
\enumo code, and could easily be written by an \enumo user,
highlighting the flexibility and power of the tool.
\begin{figure}
\begin{lstlisting}
lang = { LIT (UOP EXPR) (BOP EXPR EXPR) }
def recursive_rules($\mathcal{G}$, metric, n):
  if n == 0: return []
  rec_rules = recursive_rules($\mathcal{G}$, metric, n - 1)
  workload = iter_metric(lang, "EXPR", metric, n)
               .plug("LIT", $\mathcal{G}$.lits).plug("UOP", $\mathcal{G}$.uops).plug("BOP", $\mathcal{G}$.bops)
  rules_n = workload
              .to_egraph()
              .compress(rec_rules)
              .find_candidates()
              .minimize(rec_rules)
  return rec_rules.union(rules_n)
\end{lstlisting}
\caption{The \T{recursive_rules} function. $\mathcal{G}$ is a struct that
  specifies a grammar, containing workloads for literal expressions,
  unary operators, and binary operators. \T{metric} is the \enumo metric
  used to define an upper bound on terms, and \T{n} is the size limit.
  \T{recursive_rules} incrementally builds a ruleset by learning rules
  over terms from the domain of increasing size up to the
  specified limit.}
\label{fig:rec-rules}
\end{figure}

After running our \enumo programs,
we compared the derivability of the generated rulesets
to those produced by \ruler,
using the same grammar and interpreter.
We found that for rational arithmetic,
  \ruler learns rules over division by assuming that
  the denominator is not zero\footnote{To mitigate the resulting
  unsoundness, \ruler used a custom rule application
  strategy from the \egg~\citep{egg} library.}.
We removed this unsound
  assumption from \ruler and re-synthesized its
  rational arithmetic rules for our comparison.
We also found that for \C{bool}, \C{bv4}, and \C{bv32}, \ruler's hard-coded
  term enumeration loop does not enumerate any constant values.
As a result, \ruler fails to learn simple identities like \C{(\& ?a (~ ?a)) \rewritesto \, false}
  and \C{(<< ?a 0) \rewritesboth \, ?a}.
Failure to enumerate these constants explains why \ruler is unable to derive
  many of \enumo's rules for these domains.
We adapted the \enumo programs to not enumerate constants, and we find that
  \ruler can derive 100\% of the \C{bool} rules, 90.7\% of the \C{bv4} rules,
  and 87.8\% of the \C{bv32} rules\footnote{Both \lhs and \lhsandrhs derivability
  result in these percentages.}.

We found that \enumo rulesets were able to derive
  all of \ruler's rules using \ruler's
  own \lhsandrhs derivability metric
  (\autoref{subsec:derivability}, \autoref{table:oopsla}).
The reverse is not true.
Using the more conservative \lhs metric, \enumo rulesets derive
  a higher percentage of \ruler's rules than \ruler's rules can derive of
  \enumo's.
Both measures suggest that \enumo rulesets have greater proving power
  than their \ruler counterparts.
This result demonstrates that small (\C{<}60 LOC), simple
\enumo programs outperform \ruler while also providing users
with increased flexibility and transparency.

\subsubsection{Scaling to large grammars: A Halide case study}
\label{subsec:halideeval}
Halide~\citep{halide} is a programming language
  for high-performance
  image processing.
A major component of the Halide compiler is
  a traditional term rewriting
  system~\citep{julie-halide} that performs optimizing
  program transformations using a set of handwritten rules.
Halide has a large grammar, totaling 17 operators,
  which include boolean,
  arithmetic, and
  comparison operators.
Halide does not use an \eqsat engine
  for applying the rewrite rules.
Nevertheless, inspired by the domain,
  particularly due to the size of its grammar,
  we developed an \enumo program to evaluate
  the scalability of \enumo's workload-guided strategy.

Halide's handwritten ruleset was collected by scraping
source files from the latest commit in the Halide
  repository\footnote{\url{https://github.com/halide/Halide/commit/e7f78600e10956b44e8f214c686f310211b0d836}}
  and removing rules that we could not parse as \enumo rules.
These included rules with side conditions, rules with unsupported operators,
  and rules with unbound variables on the right-hand side of the rule.
After this process, we were left with 725 rules.

To see how prior work~\citep{ruler}
  would perform on a large domain,
  we implemented the Halide grammar in \ruler. \ruler's
  implementation was able to run for
  just one iteration (further iterations did not terminate),
  synthesizing a total of 90 rules in about 3 seconds.
  This ruleset was able to derive
  only 18 of the 725 original rules (2.5\%)
   using both derivability metrics (\lhs, \lhsandrhs).

Without leveraging guided search, \enumo scales similarly to Ruler
  for exhaustive exponential searches.
A simple \enumo program that exhaustively enumerates Halide terms up to size
  5 is able to derive 327 of Halide's 725 rules.
The exhaustive \enumo program only outperforms \ruler because \ruler
  enumerates by depth, which grows much faster than size.
An \enumo program that enumerates by depth times out after depth 2 and learns
  rules that can only derive 13 of Halide's rules, which is similar to
  \ruler's behavior.

However, the key benefit of \enumo's guided enumeration is decoupling
  the grammar and the workload size.
In \ruler, terms are enumerated exhaustively from the grammar up to a
  certain size, so with a larger grammar, \ruler hits resource limits
  faster.
In contrast, term enumeration in \enumo is separate from the
  grammar itself, allowing users to define workloads that represent
  different subsets of the search space.
With \enumo's operators, it is easy to compose workloads together,
  enabling a piecewise rather than total approach to term enumeration.
This composability makes it possible to synthesize rulesets that are
  larger and deeper than would be possible with a one-shot theory
  exploration tool like \ruler.

To evaluate whether \enumo's
  guided search would help to find deeper, more complex Halide rules,
  we wrote a 141-line \enumo program
  leveraging both exhaustive and custom enumeration.
  First, we exhaustively enumerated terms over subsets of Halide's
  operators---boolean, arithmetic,
  comparison\footnote{For both \ruler and \enumo,
  we enumerated terms over
  15 of the 17 operators, skipping $/$ and $\implies$ since
  most of those rules had side conditions.}%
  ---up to 5 atoms, beyond which point this strategy
  is computationally infeasible.
  We then enumerated
  terms with \textit{all} of
  Halide's operators up to
  4 atoms in size.
Finally, we created custom
  workloads guided by domain knowledge,
  selectively generating terms too large to be
  found using the exhaustive approach,
  such as \C{(select a (min b c) (max d c))}.
  These workloads leverage \enumo features
  such as canonicalization
  (\autoref{sec:slide}), which eliminates many duplicate terms,
  reducing one workload
  from 52,491 to 9,233 terms---an 82\%
  decrease\footnote{More details on how the Halide programs
  were written are available in \autoref{subsec:halide-appendix}.}.
 Ultimately, our \enumo program produced a ruleset of 840 rules
  capable of deriving
  84.0\% and 90.9\% of the handwritten ruleset using
  the \lhs and \lhsandrhs derivability
  metrics, respectively.

As mentioned in \autoref{subsec:derivability}, larger rulesets are not
  necessarily better than smaller rulesets.
In this case study, however, the smaller ruleset (from one iteration of \ruler)
  has measurably less proving power than the larger ruleset
  (from the \enumo program), so the additional rules are justified.
Synthesizing rulesets for large grammars is not feasible
  with tools that rely on exhaustive term enumeration,
  but with \enumo, it is possible to build
  rulesets incrementally, so grammar size is not a
  limiting factor.
This section shows that a small program
using \enumo's novel guided search finds better rules than
the state of the art can.

